\section{Introduction}
\label{sec:introduction}

Thin-film materials experiments couple solution preparation and controlled deposition with drying or curing and subsequent characterization. For films intended for corrosion protection, completing the robotic motion is not sufficient: the resulting substrates must also satisfy specified requirements for thickness uniformity, wettability, surface roughness, and corrosion performance. Automation should therefore be evaluated at both the operation level and the scientific-output level, rather than solely by whether a manipulator reaches its target.

This application poses two connected challenges. First, materials workflows span preparation and processing workstations, requiring heterogeneous robots to coordinate manipulation, inter-station transport, and container transfer. Stage completion checks and bounded recovery are needed to prevent a local failure from silently propagating through the process. Second, precision liquid operations are sensitive to perception errors, contact conditions, and differences between simulated and physical hardware. Collecting large physical training datasets can be costly, motivating simulation-based preparation followed by adaptation with a small number of real demonstrations. Whether such adaptation remains effective under held-out conditions requires separate physical testing.

We address these challenges through a hybrid robotic framework whose primary application is thin-film materials experimentation (Fig.~\ref{fig:overall_framework}). The architecture assigns preparation operations to a 7-DoF Franka Panda, inter-station delivery to a differential-drive mobile platform, and downstream coating-related operations to a 6-DoF UR3. The manipulator roles are coordinated through a common skill interface and task-level completion checks. The workflow covers solution preparation, delivery and transfer, coating, and post-processing; film characterization supplies the quality endpoints rather than being equated with motion completion.

The control architecture separates transport from precision manipulation. Constraint-aware classical planning and trajectory smoothing support mobile transport; these mechanisms target motion regularity but do not establish spill-free operation without direct liquid measurements. Posture-aware reinforcement learning supports grasping and lifting of labware. For precision operations such as pipetting and targeted dispensing, simulation demonstrations provide behavioral-cloning pretraining, domain randomization exposes the policy to variable conditions, and parameter-regularized fine-tuning uses a small set of real demonstrations to adapt the policy before physical evaluation. This offline policy adaptation is distinct from online calibration of a digital twin.

A failure-aware hierarchical state machine sequences independently implemented skills using explicit preconditions, completion guards, timeouts, and bounded retries. It is a rule-based execution mechanism, not a learned high-level reinforcement-learning policy. Timestamped synchronization of robot joints, gripper state, and object poses provides auxiliary workcell monitoring and simulation--physical trajectory comparison. Neither online parameter calibration nor twin-prediction feedback is claimed as a contribution.

The evaluation is organized around two linked questions: can the heterogeneous system reliably execute the materials workflow and produce acceptable films, and how much do domain randomization and a limited real-demonstration budget improve precision operation performance? Existing tables report individual-skill outcomes. Supplementary tests distinguish nominal conditions, held-out configurations within the training support, and conditions outside that support (Section~\ref{subsec:transfer_generalization}). The thin-film comparison considers manual operation, fixed robot trajectories, and learned control using independent substrates and material-quality measurements (Section~\ref{subsec:exp_dipcoat_quality}). Its quantitative results and complete-workflow acceptance rates remain to be supplied; they are not inferred from pooled skill success. Serial dilution is retained as an extension of the liquid-handling library, not as a coequal primary application or evidence of a completed biological assay.

The paper is organized around the following system and method elements:
\begin{itemize}
    \item \textbf{Heterogeneous robotic execution for materials experiments:} A hybrid architecture coordinates workstation manipulation and mobile transport through a reusable skill library and failure-aware task sequencing.
    \item \textbf{Few-shot sim-to-real adaptation for precision operations:} Behavioral cloning, domain randomization, and regularized real-demonstration fine-tuning form the transfer pipeline, with method and demonstration-budget comparisons.
    \item \textbf{Operation-to-material evaluation:} The evaluation separates skill success, complete-workflow execution, and film-quality acceptance; reported results are distinguished from pending quality and generalization tests.
\end{itemize}
